\documentclass[landscape]{article}
\usepackage[utf8]{inputenc}
\usepackage{graphicx}
\usepackage{tikz}
\usetikzlibrary{trees}
\usepackage{enumerate}
\usepackage{icomma}
\usepackage{amssymb}
\usepackage{eurosym}
\usepackage{tikz}
\usepackage{href-ul}
\hypersetup{
	colorlinks=true,
	linkcolor=blue,
	urlcolor=blue}
\usepackage{geometry}
\geometry{a4paper, top=15mm, left=15mm, right=15mm, bottom=15mm,
headsep=10mm, footskip=12mm}

\title{m11zugL255}

\begin{document}
\begin{tikzpicture}(10,3)
	\draw[ultra thick](2,0) --(10,0) -- (10,1.3) --(2,1.3) -- (2,0);
	\draw[fill=black](2,0)-- (0,.6) -- (2,.6) -- (2,0);
	\node at (6,.6) {\large Lösungsblatt von \href{https://www.okuyakl.de}{www.okuyakl.de}};
\end{tikzpicture}
\vspace{0.5 cm}

\noindent{\bf Aufgabe 1. a)}
$$
\renewcommand{\arraystretch}{2}
\begin{array}{|p{2 cm}|p{1.8 cm}|p{1.8 cm}|p{1.8 cm}|p{1.8 cm}|p{1.8 cm}|p{1.8 cm}|p{1.8 cm}|p{1.8 cm}|p{1.8 cm}|p{1.8 cm}|}
\hline
\textnormal{Ergebnisse} &1; 10 & 2; 11; 20 & 3; 12; 21 & 4; 13; 22 & 5; 14 ; 23 & 6; 15; 24 & 7; 16; 25 & 8; 17 & 9; 18 & 19\\
\hline
$x_i$ & 1 & 2 & 3 & 4 & 5 & 6 & 7 & 8 & 9 & 10 \\
\hline
$P(X=x_i)$ & $2\over 25$ & $3\over 25$ & $3\over 25$ &$3\over 25$ & $3\over 25$ & $3\over 25$ & $3\over 25$ & $2\over 25$ & $2\over 25$& $1\over 25$\\
\hline
\end{array}
$$
\noindent{\bf Aufgabe 1. b)}

\setlength{\unitlength}{0.5 cm}
\begin{picture}(14,5)
\put(1,1){\vector(1,0){11}}
\put(1,1){\vector(0,1){4}}

\put(0.8,0.2){0}
\put(1.8,0.2){1}
\put(2.8,0.2){2}
\put(3.8,0.2){3}
\put(4.8,0.2){4}
\put(5.8,0.2){5}
\put(6.8,0.2){6}
\put(7.8,0.2){7}
\put(8.8,0.2){8}
\put(9.8,0.2){9}
\put(10.6,0.2){10}
\put(11.5,1.2){$x_i$}

\put(-0.6,1.7){0,04 --}
\put(-0.6,2.7){0,08 --}
\put(-0.6,3.7){0,12 --}
\put(1.2,4.5){$P(X=x_i)$}

\put(2,1){\line(0,1){2}}
\put(3,1){\line(0,1){3}}
\put(4,1){\line(0,1){3}}
\put(5,1){\line(0,1){3}}
\put(6,1){\line(0,1){3}}
\put(7,1){\line(0,1){3}}
\put(8,1){\line(0,1){3}}
\put(9,1){\line(0,1){2}}
\put(10,1){\line(0,1){2}}
\put(11,1){\line(0,1){1}}
\end{picture}
\hspace{2 cm}
\begin{picture}(14,13)
\put(1,1){\vector(1,0){11}}
\put(1,1){\vector(0,1){11}}

\put(0.8,0.2){0}
\put(1.8,0.2){1}
\put(2.8,0.2){2}
\put(3.8,0.2){3}
\put(4.8,0.2){4}
\put(5.8,0.2){5}
\put(6.8,0.2){6}
\put(7.8,0.2){7}
\put(8.8,0.2){8}
\put(9.8,0.2){9}
\put(10.6,0.2){10}
\put(11.5,1.2){$x_i$}

\put(-0.4,1.7){0,1 --}
\put(-0.4,2.7){0,2 --}
\put(-0.4,3.7){0,3 --}
\put(-0.4,4.7){0,4 --}
\put(-0.4,5.7){0,5 --}
\put(-0.4,6.7){0,6 --}
\put(-0.4,7.7){0,7 --}
\put(-0.4,8.7){0,8 --}
\put(-0.4,9.7){0,9 --}
\put(-0.4,10.7){1,0 --}
\put(1.2,11.5){$P(X \le x_i)$}

\put(2,1.8){\line(1,0){1}}
\put(2,1.8){\circle*{0.3}}
\put(3,3){\line(1,0){1}}
\put(3,3){\circle*{0.3}}
\put(4,4.2){\line(1,0){1}}
\put(4,4.2){\circle*{0.3}}
\put(5,5.4){\line(1,0){1}}
\put(5,5.4){\circle*{0.3}}
\put(6,6.6){\line(1,0){1}}
\put(6,6.6){\circle*{0.3}}
\put(7,7.8){\line(1,0){1}}
\put(7,7.8){\circle*{0.3}}
\put(8,9){\line(1,0){1}}
\put(8,9){\circle*{0.3}}
\put(9,9.8){\line(1,0){1}}
\put(9,9.8){\circle*{0.3}}
\put(10,10.6){\line(1,0){1}}
\put(10,10.6){\circle*{0.3}}
\put(11,11){\line(1,0){1}}
\put(11,11){\circle*{0.3}}
\end{picture}

\noindent{\bf Aufgabe 1. c)}
$P(X=7)$ ist die Wahrscheinlichkeit, dass die Quersumme der Zahl auf der gezogenen Kugel gleich 7 ist.
Es gilt:
$$P(X=7)={3\over 25}$$
$P(X\le 4)$ ist die Wahrscheinlichkeit, dass die Quersumme der Zahl auf der gezogenen Kugel gleich 4 oder weniger ist.
Es gilt:
$$P(X\le 4)=P(X=1)+P(X=2)+P(X=3)+P(X=4)={11\over 25}$$
Diesen Wert kann man auch aus dem Diagramm der kumulativen Verteilungsfunktion ablesen.\\
$P(X < 6)$ ist die Wahrscheinlichkeit, dass die Quersumme der Zahl auf der gezogenen Kugel kleiner als 6 ist.
Es gilt:
$$P(X< 6)=P(X=1)+P(X=2)+P(X=3)+P(X=4)+P(X=5)={14\over 25}$$
Diesen Wert kann man auch aus dem Diagramm der kumulativen Verteilungsfunktion ablesen.\\
$P(5 < X \le 9)$ ist die Wahrscheinlichkeit, dass die Quersumme der Zahl auf der gezogenen Kugel größer als 5, aber höchstens 9 ist.
Es gilt:
$$P(5 < X \le 9)=P(X=6)+P(X=7)+P(X=8)+P(X=9)={10\over 25}$$

\begin{minipage}{0.3\textwidth}
\noindent{\bf Aufgabe 2. a)}
%\documentclass{article}
%\usepackage[latin1]{inputenc}
%\usepackage[utf8]{inputenc}
%\usepackage{tikz}
%\usetikzlibrary{trees}


%\begin{document}
% Set the overall layout of the tree
\tikzstyle{level 1}=[level distance=2cm, sibling distance=3.5cm]
\tikzstyle{level 2}=[level distance=2cm, sibling distance=2cm]

% Define styles for bags and leafs
\tikzstyle{bag} = [text width=4em, text centered]
\tikzstyle{end} = [circle, minimum width=3pt,fill, inner sep=0pt]

% The sloped option gives rotated edge labels. Personally
% I find sloped labels a bit difficult to read. Remove the sloped options
% to get horizontal labels. 
\begin{tikzpicture}[grow=right]
\node[bag] {Start}
    child {
        node[bag] {6}        
            child {
                node[end, label=right: {$6 \qquad P(66)=\frac{2}{7}\cdot\frac{2}{7}$}]{}
                edge from parent
                node[above] {}
                node[below]  {$\frac{2}{7}$}
            }
            child {
                node[end, label=right:
                    {$\bar{6}\qquad P(6\bar{6})=\frac{2}{7}\cdot\frac{5}{7}$}] {}
                edge from parent
                node[above] {}
                node[below]  {$\frac{5}{7}$}
            }
            edge from parent 
            node[above] {}
            node[below]  {$\frac{2}{7}$}
    }
    child {
        node[bag] {$\bar{6}$}        
        child {
                node[end, label=right:
                    {$6 \qquad P(\bar{6}6)=\frac{5}{7}\cdot\frac{2}{7}$}] {}
                edge from parent
                node[above] {}
                node[below]  {$\frac{2}{7}$}
            }
            child {
                node[end, label=right:
                    {$\bar{6} \qquad P(\bar{6}\bar{6})=\frac{5}{7}\cdot\frac{5}{7}$}] {}
                edge from parent
                node[above] {}
                node[below]  {$\frac{5}{7}$}
            }
        edge from parent         
            node[above] {}
            node[below]  {$\frac{5}{7}$}
    };
\end{tikzpicture}
%\end{document}
$$
\renewcommand{\arraystretch}{2}
\begin{array}{|c|c|c|c|}
\hline
x_i & 0 & 1 & 2 \\
\hline
P(X=x_i) & {25\over 49}& {20 \over 49} & {4 \over 49} \\
\hline
\end{array}
$$
\end{minipage}
\begin{minipage}{0.3\textwidth}

{\bf Aufgabe 2. b)}
$$
\renewcommand{\arraystretch}{2}
\begin{array}{|c|c|c|c|c|c|c|}
\hline
x_i & 0 & 1 & 2 & 3 & 4 & 6 \\
\hline
P(X=x_i) & {1\over 9} & {2\over 9}& {1\over 9}& {1\over 9}& {2\over 9}& {1\over 9}\\
\hline
\end{array}
$$

{\bf Aufgabe 2. c)}
$$
\renewcommand{\arraystretch}{2}
\begin{array}{|c|c|c|c|}
\hline
x_i & -2,50 & 1 & 2 \\
\hline
P(X=x_i) & 0,4 & 0,4 & {1 \over 5} \\
\hline
\end{array}
$$
\end{minipage}
\begin{minipage}{0.3\textwidth}
\noindent{\bf Aufgabe 2. d)}\\
\includegraphics[width=6 cm]{baum2d}
$$
\renewcommand{\arraystretch}{2}
\begin{array}{|c|c|c|c|c|}
\hline
x_i & 0 & 1 & 2 & 3\\
\hline
P(X=x_i) & 0,0067 & 0,0576 & 0,254 & 0,681\\
\hline
\end{array}
$$
\end{minipage}

\begin{center}
	\includegraphics[width=7 cm]{../../viecher/endcomic.pdf}
	
	Hier geht es zurück zum \href{https://www.okuyakl.de/math/m11zugL255/aa255.pdf}{Aufgabenblatt}
\end{center}

\end{document}
